\documentclass[a4paper, 12pt]{article}
\usepackage[serbian]{babel}
\usepackage[a4paper, margin=2.5cm]{geometry} 
\usepackage{fancyhdr} 
\usepackage{titlesec} 
\usepackage{tocloft} 
\usepackage{biblatex} 
\usepackage{csquotes} 
\usepackage[T1]{fontenc}
\usepackage{hyperref}

\hypersetup{
    colorlinks = true,
    linkcolor = black,
    filecolor = green,      
    urlcolor = blue,
    pdftitle = {Pregled gradiva kursa},
    pdfpagemode = FullScreen,
    }

\setlength{\headheight}{14.5pt} % Fix fancyhdr warning

\pagestyle{fancy}
\fancyhf{} % Clear default header/footer
\fancyhead[L]{\small \leftmark} % Header left
\fancyfoot[C]{\thepage} % Page number in center footer

\titleformat{\section}{\large\bfseries}{\thesection}{1em}{}


\title{Pregled gradiva kursa \\ \large Računarstvo i društvo}
\author{Matija Stanković}
\date{Januar 2026}

\renewcommand{\contentsname}{Sadržaj}

\begin{document}

\maketitle
\newpage
\tableofcontents
\newpage


\section{Istorijski pregled razvoja računarstva}
\subsection{Pregled gradiva}
Razvoj računarstva predstavlja dugotrajan proces u kome su se istovremeno razvijali
hardver, softver i sistemi za prenos informacija. Pretečama današnjih računara smatraju se elektronse računske mašine, koje su omogućile automatsku obradu
velikog broja podataka. Među najznačajnijim ranim računarima izdvajaju se uređaji poput
ABC-a i ENIAC-a, koji su najsličniji današnjim računarima opšte namene.

Razvoj računara može se podeliti u četiri faze, u zavisnosti od tehnologije koja je u
određenom periodu bila dominantna:
\begin{itemize}
    \item Prva faza – \textbf{računari zasnovani na vakuumskim cevima}
    \item Druga faza – \textbf{računari zasnovani na tranzistorima}
    \item Treća faza – \textbf{računari zasnovani na integrisanim kolima}
    \item Četvrta faza – \textbf{mikroprocesorski i personalni računari}
\end{itemize}

U prvoj fazi razvoja računari su koristili vakuumske cevi, zbog čega su bili velikih
dimenzija, imali veliku potrošnju energije i bili skloni čestim kvarovima. I pored
toga, ovi računari su predstavljali značajan napredak jer su omogućili znatno bržu
obradu podataka u odnosu na ranije mehaničke sisteme.

Početak druge faye razvoja vezuje se za uvođenje tranzistora, koji su zamenili vakuumske cevi.
Na ovaj način povećana je pouzdanost računara, smanjena njihova veličina i potrošnja
energije, a performanse su dodatno unapređene.

Treća faza razvoja obeležena je upotrebom integrisanih kola, koja su omogućila smeštanje
velikog broja elektronskih komponenti na mali prostor. Ovo je dovelo do daljeg smanjenja
računara i njihove šire primene u različitim oblastima.

Razvoj mikroprocesora označava početak četvrte faze, u kojoj dolazi do pojave personalnih
računara. Računari postaju dostupni širem krugu korisnika i pronalaze primenu u svakodnevnom
životu, obrazovanju i poslovanju.

Paralelno sa razvojem računara javila se potreba za njihovim međusobnim povezivanjem.
Razvoj računarskih mreža započinje projektom ARPANET, koji se smatra pretečom današnjeg
interneta. 

Savremeno računarstvo karakteriše brz razvoj novih tehnologija, kao što su
internet, mobilni uređaji, veštačka inteligencija, virtuelna realnost i blockchain, koje
imaju značajan uticaj na savremeno društvo.

\subsection{Predložene teme}
\begin{itemize}
    \item Upotreba informacionih tehnologija u nastavi – stavovi i mišljenja nastavnika i učenika
    \begin{itemize}
        \item \url{https://nardus.mpn.gov.rs/bitstream/id/68509/Disertacija.pdf}
    \end{itemize}
    \item Programiranje za osnovce
    \begin{itemize}
        \item \url{https://www.sk.rs/2015/10/skpr01.html}
    \end{itemize}
    
\end{itemize}

\section{Etika}
\subsection{Pregled gradiva}
\textbf{Etika} je grana filozofije koja se bavi proučavanjem moralnih vrednosti, normi i pravila
koja usmeravaju ljudsko ponašanje, odnosno razlikovanjem ispravnog i pogrešnog
postupanja.

Etika se deli na sledeće podoblasti:
\begin{itemize}
    \item \textbf{Metaetika} - proučava značenje moralnih pojmova i pitanje šta zapravo
    znači kada kažemo da je nešto dobro ili loše.
    \item \textbf{Normativna etika} - bavi se pravilima i principima koji određuju kako bi
    ljudi trebalo da se ponašaju.
    \item \textbf{Primenjena etika} - razmatra konkretne etičke probleme i dileme u
    stvarnim životnim situacijama.
\end{itemize}

Etička teorija predstavlja skup principa i pravila koji se koriste za procenu
moralnosti postupaka i donošenje etičkih odluka.

Etičke teorije mogu se podeliti na:
\begin{itemize}
    \item \textbf{Subjektivni i kulturni relativizam} - moralni sudovi zavise od
    pojedinca ili kulture i ne postoje univerzalna moralna pravila.
    \item \textbf{Etički egoizam} - moralno ispravno je ono što je u ličnom interesu
    pojedinca
    \item \textbf{Kantijanizam} - postupci se vrednuju prema dužnosti i univerzalnosti
    pravila, a ne prema posledicama
    \item \textbf{Utilitarizam} - moralnost postupka se procenjuje prema količini koristi
    koju donosi najvećem broju ljudi.
    \item \textbf{Teorija društvenog ugovora} - moralna pravila nastaju dogovorom članova
    društva radi očuvanja reda i pravednosti.
\end{itemize}

Etika tehnologije predstavlja deo primenjene etike koji proučava moralne
posledice razvoja i primene tehnologije u savremenom društvu.

Računarska etika odnosi se na etička pitanja vezana za upotrebu računara i
informacionih tehnologija, kao što su privatnost, bezbednost, odgovornost i
intelektualna svojina.

U okviru računarske etike izdvajaju se sledeći etički stavovi i pristupi:
\begin{itemize}
    \item \textbf{Profesionalni pristup} - naglašava poštovanje profesionalnih kodeksa i
    odgovornost stručnjaka.
    \item \textbf{Akademski pristup} - usmeren na teorijsko proučavanje etičkih problema i
    obrazovanje u oblasti računarstva.
    \item \textbf{Inovativni pristup} - razmatra nove etičke izazove nastale usled brzog
    razvoja informacionih tehnologija.
\end{itemize}

\newpage
Deset zapovesti računarske etike:
\begin{itemize}
    \item Ne koristiti računar kako bi se nanosila šteta drugim ljudima.
    \item Ne ometati rad drugih ljudi koristeći računar.
    \item Ne pregledati, menjati ili uništavati tuđe fajlove bez dozvole.
    \item Ne koristiti računar za krađu ili prevaru.
    \item Ne koristiti računar za širenje lažnih ili obmanjujućih informacija.
    \item Ne koristiti softver ili druge digitalne resurse bez odgovarajuće licence
    ili dozvole.
    \item Ne prisvajati tuđi intelektualni rad bez navođenja izvora.
    \item Razmišljati o društvenim posledicama programa i sistema koje razvijamo ili
    koristimo.
    \item Koristiti računar uz poštovanje privatnosti i dostojanstva drugih ljudi.
    \item Uvek koristiti računar na način koji pokazuje poštovanje i odgovornost
    prema zajednici.
\end{itemize}

\subsection{Predložene teme}
\begin{itemize}
    \item Etika i integritet u javnom sektoru
    \begin{itemize}
        \item \url{https://www.youtube.com/playlist?list=PLHBowDEf-zianeStj6g72KC8KGT85jKar}
    \end{itemize}
    \item Radna etika
    \begin{itemize}
        \item \url{https://proeduca.net/sta-sve-cini-radnu-etiku/}
    \end{itemize}
    
\end{itemize}

\section{Upotreba interneta, spam poruke}
\subsection{Pregled gradiva}
\textbf{Spam poruke} predstavljaju neželjene poruke koje se masovno šalju velikom broju
korisnika bez njihove saglasnosti. Najčešće se pojavljuju u okviru elektronske
pošte. Osnovna karakteristika spam poruka jeste da primalac nije inicirao komunikaciju.

Prve spam poruke pojavile su se sa razvojem elektronske pošte i širenjem interneta.
Masovno slanje poruka postalo je moguće zbog jednostavne i jeftine distribucije
elektronske pošte, što je dovelo do pojave prvih zloupotreba ovog vida komunikacije.

Osnovna prednost spam poruka za pošiljaoce ogleda se u sledećem:
\begin{itemize}
    \item mogućnost masovnog slanja poruka,
    \item veoma niski troškovi reklamiranja,
    \item potencijalna dobit čak i uz mali broj odgovora.
\end{itemize}

Sadržaj spam poruka najčešće obuhvata:
\begin{itemize}
    \item reklamne i promotivne poruke,
    \item lažne ponude i obmanjujuće informacije,
    \item poruke sa linkovima ka sumnjivim internet stranicama,
    \item priloge koji mogu predstavljati bezbednosni rizik.
\end{itemize}

Spam poruke predstavljaju značajan problem jer:
\begin{itemize}
    \item opterećuju komunikacione i mrežne resurse,
    \item otežavaju normalnu elektronsku komunikaciju (gubitak novca za kompanije),
    \item smanjuju propusnu moć Interneta.
\end{itemize}

Resursi za slanje spam poruka obezbeđuju se korišćenjem automatizovanih programa,
zloupotrebom računara korisnika i formiranjem mreža zaraženih računara, koje
omogućavaju masovno slanje poruka bez znanja vlasnika sistema.

Radi smanjenja količine spam poruka razvijeni su spam filteri. Njihova uloga je da
prepoznaju i izdvoje neželjene poruke na osnovu sadržaja poruke, podataka o
pošiljaocu i obrasca ponašanja, pre nego što poruka stigne do korisnika.

\vspace{\baselineskip}

\textbf{World Wide Web} (www) predstavlja servis interneta koji omogućava pristup dokumentima koji su međusobno povezani linkovima. 

Glavne karakteristike World Wide Web-a su:
\begin{itemize}
    \item decentralizovan sistem,
    \item svaki resurs (dokument) ima jedinstvenu adresu,
    \item zasnovan je na internetu.
\end{itemize}

Internet i World Wide Web ne predstavljaju isto. Internet je globalna računarska mreža
, a World Wide Web jedan od servisa koji funkcioniše na Internetu i koristi njegovu
infrastrukturu za pristup dokumentima i informacijama.

Internet se koristi za:
\begin{itemize}
    \item razmenu elektronske pošte,
    \item pristup World Wide Web servisima,
    \item razmenu datoteka,
    \item komunikaciju u realnom vremenu,
    \item obrazovanje i informisanje,
    \item zabavu i društvene mreže.
\end{itemize}
\subsection{Predložene teme}

\begin{itemize}
    \item Kako prepoznati i sprečiti spam poruke
    \begin{itemize}
        \item \url{https://www.plus.rs/blog/kako-prepoznati-i-spreciti-spam-poruke/}
    \end{itemize}
    \item Mediji VS društvene mreže u Srbiji
    \begin{itemize}
        \item \url{https://www.bbc.com/serbian/lat/srbija-54718120}
    \end{itemize}
    
\end{itemize}


\section{Cenzura i sloboda govora, deca i neprikladan sadržaj}
\subsection{Pregled gradiva}
\textbf{Cenzura} predstavlja pokušaj da se suzbije javni pristup informacijama ili
materijalima koji se smatraju uvredljivim ili štetnim. 
Kroz istoriju cenzuru su najčešće sprovodile državne i religiozne institucije, a njen značaj 
se povećava razvojem tehnologije i sredstava komunikacije.

Vrste cenzure:
\begin{itemize}
    \item \textbf{Direktna cenzura} - ima za cilj da spreči pojavljivanje neprihvatljivog
    materijala ili da ograniči njegovu dostupnost javnosti.
    \item \textbf{Samocenzura} - predstavlja kontrolu sopstvenog govora i ponašanja radi
    izbegavanja mogućih posledica.
\end{itemize}

Direktna cenzura se sastoji iz tri osnovna oblika:
\begin{itemize}
    \item \textbf{Monopol od strane vladinih institucija} - država poseduje i kontroliše
    medije i izvore informacija i tako u potpunosti upravlja tokom informacija.
    \item \textbf{Pregled publikacije} - kontrola i zabrana objavljivanja sadržaja koji
    se smatra štetnim, uvredljivim ili opasnim po državu.
    \item \textbf{Licenciranje i registracija} - ograničavanje pristupa medijima kroz
    dodelu dozvola i licenci, pri čemu vlast odlučuje ko može da emituje sadržaj.
\end{itemize}

Samocenzura na društvenim mrežama ogleda se u tome da korisnici sami odlučuju
koji sadržaj će podeliti, sa kim i u kojoj meri.

Teškoće cenzurisanja na Internetu proizlaze iz njegove globalne prirode,
neprekidnog širenja i nemogućnosti potpune kontrole objava van granica jedne
države.

Manipulisanje Internetom od strane države podrazumeva ograničavanje pristupa
određenim sadržajima ili servisima, pri čemu se standardi prihvatljivosti
razlikuju od države do države.


\vspace{\baselineskip}

\textbf{Sloboda govora} predstavlja pravo pojedinca da slobodno izražava mišljenja bez straha da će u tome biti sprečen ili da će biti kažnjen za izraženo mišljenje.

Sloboda govora danas suočava se sa brojnim izazovima, jer se pod izgovorom
političke korektnosti i zaštite od uvreda uvode nova ograničenja izražavanja, naročito u digitalnom okruženju.

\vspace{\baselineskip}

\textbf{Deca i neprikladan sadržaj} predstavljaju poseban problem na Internetu, jer su
deca često izložena pornografskom, nasilnom ili uznemirujućem sadržaju, naročito
sa širenjem pametnih telefona i stalnim pristupom mreži.

Prečišćivač Veba je alat koji filtrira i onemogućava prikaz nepoželjnih
internet stranica tokom njihovog učitavanja u pregledaču.

Prečišćivače Veba koriste:
\begin{itemize}
    \item roditelji - radi zaštite dece od neprikladnog sadržaja.
    \item poslovne grupe - kako bi ograničile pristup sadržajima nevezanim za posao.
    \item obrazovne ustanove - u cilju sprečavanja uvredljivog i ometajućeg sadržaja.
\end{itemize}

\vspace{\baselineskip}

\textbf{Sexting} predstavlja slanje poruka ili elektronske pošte koje sadrže eksplicitan
sadržaj, a najčešće je prisutan među adolescentima.

Sexting se deli na sledeće tipove:
\begin{itemize}
    \item razmena slika između partnera,
    \item razmena slika između partnera koje su prosleđene trećim licima,
    \item razmena sadržaja između osoba koje još nisu u vezi, ali postoji očekivanje
    budućeg odnosa.
\end{itemize}

\subsection{Predložene teme}

\begin{itemize}
    \item Sloboda govora u digitalnom prostoru
    \begin{itemize}
        \item \url{https://insajder.net/vesti/sloboda-govora-u-digitalnom-prostoru-ko-stoji-iza-gasenja-instagram-naloga-medija-i-nvo}
    \end{itemize}
    \item Zloupotreba dece na internetu
    \begin{itemize}
        \item \url{https://kliknibezbedno.wordpress.com/zloupotrebe-dece-na-internetu/}
    \end{itemize}
    
\end{itemize}

\section{Internet prevare i zavisnost od interneta}
\subsection{Pregled gradiva}
\textbf{Internet prevara} predstavlja oblik obmane u digitalnom okruženju u kome napadač
koristi Internet kako bi doveo korisnika u zabludu i ostvario
ličnu ili finansijsku korist. Ovakve prevare se često oslanjaju na nepažnju korisnika,
nedovoljno znanje ili poverenje u digitalne servise.

Internet prevare mogu se podeliti na sledeće oblike:
\begin{itemize}
    \item \textbf{Krađa identiteta},
    \item \textbf{Phishing},
    \item \textbf{Nigerijska prevara},
    \item \textbf{Lažne ocene proizvoda i usluga}.
\end{itemize}

Krađa identiteta je zloupotreba identiteta druge osobe, odnosno njenih ličnih podataka, sa
ciljem ostvarivanja prava koja su dozvoljena samo vlasniku identiteta.

Phishing predstavlja najčešći oblik Internet prevare. Napadač se predstavlja lažno kao institucija ili pojedinac od poverenja i zahteva od korisnika unos ličnih podataka koje kasnije zloupotrebljava.

Mediji koji se najčešće koriste za phishing:
\begin{itemize}
    \item \textbf{Elektronska pošta} - lažni emailovi koji zahtevaju unos ličnih podataka,
    \item \textbf{Lažne internet stranice} - kopije legitimnih sajtova,
    \item \textbf{Mobilni telefoni} - SMS poruke ili pozivi sa obmanjujućim sadržajem.
\end{itemize}

Nigerijska prevara predstavlja oblik prevare u kome napadač korisniku nudi
nepostojeću finansijsku dobit ili nasledstvo, uz zahtev da se unapred pošalje novac
ili lični podaci kako bi se transakcija navodno realizovala.

Lažne ocene proizvoda i usluga predstavljaju oblik obmane u kome se veštački
povećava ili smanjuje rejting proizvoda, usluga ili firmi (najčešće kroz neki vid ocenjivanja istih na Internetu), kako bi se uticalo na odluke korisnika prilikom kupovine ili izbora.

Netačne informacije predstavljaju informacije koje nisu proverene, nisu tačne ili
su namerno plasirane sa ciljem manipulacije korisnicima. Ovo može biti dosta opasno, jer su informacije na Internetu lako dostupne svima, i iakoInternet može biti sjajno mesto za informisanje moramo naučiti da kritički rasudjujemo o informacijama koje pokupimo sa Interneta.

Google koristi algoritam za procenu pouzdanosti Internet stranica koji se zasniva na broju i kvalitetu linkova koji vode ka određenoj stranici. Ovaj mehanizam funkcioniše po principu „glasanja“, pri čemu svaka stranica koja upućuje link predstavlja jedan glas. Glasovi stranica koje se smatraju pouzdanim imaju veću težinu u odnosu na glasove nepoznatih ili nepouzdanih stranica, čime se obezbeđuje relevantniji i kvalitetniji prikaz rezultata pretrage.

Vikipedija predstavlja otvorenu enciklopediju čiji sadržaj može uređivati veliki
broj korisnika, zbog čega može sadržati netačne ili nepotpune informacije, ali se
oslanja na sistem provera i ispravki kako bi se greške vremenom uklonile.

\vspace{\baselineskip}

\textbf{Zavisnost od Interneta} je vrlo teško definisati, a još uvek postoje podeljena mišljenja oko toga da li zavisnost od Interneta uopšte postoji.
U nekim zemljama zavisnost od Interneta karakteriše se kao psihički poremećaj.

Neki od pokazatelja zavisnostu od Interneta su:
\begin{itemize}
    \item gubitak kontrole nad vremenom provedenim na Internetu,
    \item zanemarivanje obaveza i svakodnevnih aktivnosti,
    \item osećaj nelagodnosti ili nervoze bez pristupa Internetu,
    \item potreba za sve dužim korišćenjem interneta radi zadovoljstva,
    \item konstantno razmišljanje o korišćenju Interneta
\end{itemize}

Zavisnost od Interneta može izazvati brojne probleme, među kojima su depresija, anksioznost, promene raspoloženja, osećaj krivice, usamljenost, glavobolja, gojaznost...

\subsection{Predložene teme}

\begin{itemize}
    \item Lečenje zavisnosti od Interneta
    \begin{itemize}
        \item \url{https://drvorobjev.com/zavisnost-od-interneta-izlecenje-na-klinici-dr-vorobjev/}
    \end{itemize}
    \item Internet prevare i penzioneri
    \begin{itemize}
        \item \url{https://www.rts.rs/vesti/drustvo/5668662/nova-internet-prevara-penzionera-lazni-profili-nude-udvostrucivanje-penzije---kako-se-zastititi.html}
    \end{itemize}
    
\end{itemize}


\section{Bezbednost interneta}
\subsection{Pregled gradiva}
\textbf{Hakovanje} se ranije vezivalo pre svega za radoznalost i tehničko istraživanje
računarskih sistema, dok se danas najčešće povezuje sa neovlašćenim pristupom
računarima i mrežama u cilju krađe podataka ili ostvarivanja materijalne koristi.
Haker je osoba koja neovlašćeno pristupa računarskim sistemima i mrežama.

Korisničko ime i šifra predstavljaju osnovni mehanizam autentifikacije korisnika,
zbog čega su česta meta napada. Slabe lozinke, kratke kombinacije i ponovna upotreba
iste šifre značajno povećavaju rizik od kompromitovanja naloga.

Najčešće korišćene tehnike hakovanja su:
\begin{itemize}
    \item \textbf{Metod grube sile} - isprobavanje velikog broja kombinacija lozinki,
    \item \textbf{Metod rečnika} - korišćenje liste najčešćih lozinki.
\end{itemize}

Napad u okviru besplatne Wi-Fi mreže zasniva se na presretanju saobraćaja između
korisnika i veb stranice, pri čemu je HTTP komunikacija posebno ranjiva u odnosu
na HTTPS.

Firesheep je ekstenzija pomoću koje svaki korisnik može da preuzme kolačiće neenriptovane veb strane.
Upotreba alata poput Firesheep-a postavlja ozbiljna etička pitanja, jer omogućava
zloupotrebu tehničkog znanja i narušavanje privatnosti korisnika.

Zakonske regulative predviđaju novčane kazne i kazne zatvora za neovlašćeno hakovanje
računarskih sistema i mreža.

\textbf{Zlonamerni softver} je softver namenjen nanošenju štete korisnicima i sistemima.

Neki od oblika zlonamernih softvera su:
\begin{itemize}
\item \textbf{Virus} - zlonamerni program koji se vezuje za drugi program i širi
    se njegovim pokretanjem, a prenosi se razmenom fajlova i spoljašnjih medija.

    \item \textbf{Internet crv} - samostalni program koji se širi mrežom korišćenjem
    bezbednosnih propusta bez potrebe za korisničkom interakcijom.

    \item \textbf{Crvi u aplikacijama za instant poruke} - šire se slanjem zlonamernih
    linkova ili fajlova kontaktima korisnika.

    \item \textbf{Conficker} - zlonamerni crv koji je formirao veliku mrežu zaraženih
    računara i omogućavao dalju distribuciju malvera.

    \item \textbf{Ukrštanje veb lokacija (Cross-Site Scripting)} - napad u kome se
    zlonamerna skripta ubacuje u veb stranicu radi krađe korisničkih podataka.

    \item \textbf{Usputno preuzimanje podataka (Drive-By Downloads)} - automatsko
    preuzimanje zlonamernog softvera prilikom posete kompromitovanoj veb stranici.

    \item \textbf{Trojanski konj} - zlonamerni softver koji se predstavlja kao koristan
    program, dok u pozadini izvršava štetne aktivnosti.

    \item \textbf{Spyware} - softver koji potajno prikuplja informacije o korisniku i
    prosleđuje ih trećim licima.

    \item \textbf{Adware} - softver koji prikazuje neželjene reklame i često prati
    aktivnosti korisnika.

    \item \textbf{Kolačići za praćenje} - služe za evidentiranje ponašanja korisnika na
    internetu i mogu predstavljati rizik po privatnost.

    \item \textbf{Posmatrači sistema (system monitors)} - prate aktivnosti korisnika i
    beleže podatke o radu sistema.

    \item \textbf{Mreža botova} - grupa zaraženih računara koji se centralno
    kontrolišu i koriste za masovne napade ili širenje zlonamernog softvera.
\end{itemize}
Osnovne mere zaštite od zlonamernog softvera obuhvataju:
\begin{itemize}
    \item korišćenje jakih i jedinstvenih lozinki,
    \item redovno ažuriranje operativnog sistema i aplikacija,
    \item upotrebu antivirusnog softvera,
    \item korišćenje firewall-a,
    \item oprezno ponašanje na internetu.
\end{itemize}

\subsection{Predložene teme}

\begin{itemize}
    \item Slučaj Kevina Mitnika
    \begin{itemize}
        \item \url{https://www.ebsco.com/research-starters/history/arrest-hacker-kevin-mitnick}
    \end{itemize}
    \item Hakovanje Državne Lutrije Srbije
    \begin{itemize}
        \item \url{https://www.youtube.com/watch?v=ZLQ4sigLbTE}
    \end{itemize}
    
\end{itemize}


\section{Privatnost na internetu}
\subsection{Pregled gradiva}

U martu 2018. godine bivši zaposleni firme Kembridž Analitika otkrio je javnosti da je ova firma kupila od društvene mreže Fejsbuk privatne podatke bar 50 miliona korisnika bez njihove saglasnosti. 

\textbf{Privatnost} predstavlja pravo pojedinca da kontroliše koje informacije o sebi
deli sa drugima i na koji način se te informacije koriste.

Privatnost na Internetu odnosi se na zaštitu ličnih podataka korisnika u digitalnom
okruženju. Percepcija privatnosti se razlikuje od korisnika do korisnika i
često zavisi od nivoa informisanosti i poverenja u digitalne platforme.

Privatnost i društvene mreže su usko povezane, jer korisnici često dobrovoljno
objavljuju lične podatke, fotografije i informacije o svom životu, čime smanjuju
nivo sopstvene privatnosti.

\vspace{\baselineskip}

\textbf{Kolačići} predstavljaju male tekstualne fajlove koje veb sajtovi čuvaju na uređaju
korisnika radi pamćenja podešavanja (na primer ime, preyime, adresa, postanski broj itd. pri online kupovini).

Kolačići se mogu podeliti na:
\begin{itemize}
    \item \textbf{Sesijske kolačiće} - privremeni kolačići koji se brišu nakon zatvaranja
    pregledača,
    \item \textbf{Trajne kolačiće} - ostaju sačuvani na uređaju korisnika duži vremenski
    period i koriste se za praćenje i pamćenje korisničkih podešavanja. Dele se na:
    \begin{itemize}
        \item \textbf{Kolačiće prve strane} - postavlja ih veb sajt koji korisnik posećuje i  najčešće se koriste za poboljšanje korisničkog iskustva,
        \item \textbf{Kolačiće treće strane} - postavljaju ih spoljašnji servisi, najčešće
        u svrhu dobijanja informacija o tome kada korinik poseti sajt na kome oni imaju svoje kolačiće.
    \end{itemize}
\end{itemize}


Zakonom o kolačićima regulisano je da korisnici moraju biti obavešteni o korišćenju
kolačića i dati saglasnost pre njihovog skladištenja, osim u slučajevima kada su
kolačići neophodni za funkcionisanje sajta.

Informacije koje se prikupljaju putem kolačića uključuju:
\begin{itemize}
    \item podatke o posećenim stranicama,
    \item vreme zadržavanja na sajtu,
    \item korisničke preferencije i podešavanja,
    \item podatke o uređaju i pregledaču.
\end{itemize}

Zaštita od kolačića može se ostvariti:
\begin{itemize}
    \item podešavanjem opcija privatnosti u pregledaču,
    \item brisanjem ili blokiranjem kolačića,
    \item korišćenjem režima privatnog pregleda (incognito).
\end{itemize}

\textbf{Opšta zaštita privatnosti na Internetu} podrazumeva odgovorno ponašanje korisnika,
ograničavanje deljenja ličnih podataka, korišćenje bezbednih lozinki, enkripcije i
pažljivo upravljanje podešavanjima privatnosti na Internetu.


\subsection{Predložene teme}

\begin{itemize}
    \item Kako sačuvati privatnost na Internetu
    \begin{itemize}
        \item \url{https://digitalni-vodic.ucpd.rs/zastita-licnih-podataka-i-privatnosti-na-internetu/?lng=lat}
    \end{itemize}
    \item Praćenje aktivnosti na Internetu
    \begin{itemize}
        \item \url{https://www.informacija.rs/Vesti/Ko-vas-i-kako-prati-na-internetu.html}
    \end{itemize}
    
\end{itemize}

\end{document}
